\section{Thực khuẩn thể và phân loại lối sống}
\label{sec:bacteriophage}

\subsection{Giới thiệu về thực khuẩn thể}

Thực khuẩn thể (bacteriophage), hay còn gọi tắt là phage, là các virus có khả năng xâm nhiễm và nhân lên bên trong tế bào vi khuẩn. Thuật ngữ "bacteriophage" có nguồn gốc từ tiếng Hy Lạp, trong đó "bacteria" có nghĩa là vi khuẩn và "phagein" có nghĩa là ăn, phản ánh khả năng tiêu diệt vi khuẩn của chúng. Thực khuẩn thể được phát hiện độc lập bởi hai nhà khoa học: Frederick Twort tại Anh năm 1915 và Félix d'Hérelle tại Pháp năm 1917, với d'Hérelle là người đặt tên "bacteriophage" cho các tác nhân này.

Thực khuẩn thể là dạng sống phong phú nhất trên Trái Đất, với số lượng ước tính lên tới $10^{31}$ cá thể \cite{suttle2007marine}. Chúng có mặt ở hầu hết mọi môi trường có vi khuẩn, từ đại dương, đất, không khí, đến cơ thể động vật và thực vật. Trong môi trường biển, mật độ phage có thể đạt tới $10^7$ virion trên mililít nước, vượt xa số lượng vi khuẩn. Sự phong phú này phản ánh vai trò sinh thái quan trọng của phage trong việc điều hòa quần thể vi khuẩn và duy trì cân bằng sinh thái vi sinh vật.

\subsection{Cấu trúc và thành phần của thực khuẩn thể}

Về mặt cấu trúc, thực khuẩn thể bao gồm vật liệu di truyền (DNA hoặc RNA) được bao bọc bởi một lớp vỏ protein gọi là capsid. Capsid bảo vệ vật liệu di truyền khỏi các yếu tố môi trường và đóng vai trò trong quá trình nhận diện và gắn kết với tế bào vi khuẩn chủ. Phage có thể có nhiều hình dạng khác nhau, nhưng cấu trúc phổ biến nhất là dạng đầu-đuôi (head-tail morphology), bao gồm:

\begin{itemize}
    \item \textbf{Đầu (head/capsid):} Chứa vật liệu di truyền, thường có hình dạng icosahedral (20 mặt)
    \item \textbf{Đuôi (tail):} Cấu trúc ống rỗng giúp phage gắn vào bề mặt vi khuẩn và tiêm vật liệu di truyền vào tế bào chủ
    \item \textbf{Sợi đuôi (tail fibers):} Các cấu trúc protein giúp nhận diện và gắn kết đặc hiệu với các thụ thể trên bề mặt vi khuẩn chủ
    \item \textbf{Đế đuôi (baseplate):} Cấu trúc phức tạp nối giữa đầu và đuôi, đóng vai trò trong việc xuyên thủng màng tế bào vi khuẩn
\end{itemize}

Bộ gen của phage có thể là DNA sợi đơn (ssDNA), DNA sợi đôi (dsDNA), RNA sợi đơn (ssRNA), hoặc RNA sợi đôi (dsRNA), với dsDNA là dạng phổ biến nhất. Kích thước bộ gen phage biến thiên rộng, từ vài nghìn base pairs (kb) đến hơn 500 kb, mã hóa từ vài gen đến hàng trăm gen.

\subsection{Chu trình sống của thực khuẩn thể}

Thực khuẩn thể được phân loại thành hai nhóm chính dựa trên chiến lược nhân lên: thực khuẩn thể độc lực (virulent phage) và thực khuẩn thể ôn hòa (temperate phage).

\subsubsection{Chu trình phân giải (Lytic cycle)}

Thực khuẩn thể độc lực chỉ thực hiện chu trình phân giải, bao gồm các bước sau:

\begin{enumerate}
    \item \textbf{Hấp phụ (Adsorption):} Phage gắn kết đặc hiệu với các thụ thể trên bề mặt tế bào vi khuẩn chủ thông qua các protein nhận diện trên sợi đuôi

    \item \textbf{Tiêm nhiễm (Injection):} Phage tiêm vật liệu di truyền của mình vào tế bào vi khuẩn, trong khi capsid protein thường ở lại bên ngoài

    \item \textbf{Tổng hợp sinh học (Biosynthesis):} DNA của phage chiếm đoạt bộ máy tổng hợp của tế bào chủ, sử dụng ribosome, enzyme, và các thành phần tế bào để sao chép DNA phage và tổng hợp các protein cấu trúc của phage

    \item \textbf{Lắp ráp (Assembly):} Các thành phần protein và DNA phage mới được tổng hợp tự lắp ráp thành các virion hoàn chỉnh bên trong tế bào vi khuẩn

    \item \textbf{Phân giải (Lysis):} Phage tổng hợp các enzyme phân giải (như lysozyme) để phá vỡ thành tế bào vi khuẩn, giải phóng hàng chục đến hàng trăm virion con để xâm nhiễm các tế bào vi khuẩn khác
\end{enumerate}

Chu trình phân giải thường hoàn thành trong vòng 20-40 phút, tùy thuộc vào loại phage và điều kiện môi trường.

\subsubsection{Chu trình lysogen (Lysogenic cycle)}

Thực khuẩn thể ôn hòa có khả năng thực hiện cả chu trình phân giải và chu trình lysogen. Trong chu trình lysogen:

\begin{enumerate}
    \item \textbf{Tích hợp (Integration):} Sau khi tiêm vào tế bào chủ, DNA của phage được tích hợp vào nhiễm sắc thể của vi khuẩn thông qua quá trình tái tổ hợp đặc hiệu vị trí (site-specific recombination), được xúc tác bởi enzyme integrase

    \item \textbf{Trạng thái prophage:} DNA phage được tích hợp (gọi là prophage) tồn tại ở trạng thái không hoạt động, được nhân lên cùng với bộ gen vi khuẩn qua nhiều thế hệ. Prophage được duy trì ở trạng thái im lặng nhờ protein repressor ngăn chặn biểu hiện các gen phân giải

    \item \textbf{Cảm ứng (Induction):} Khi gặp các điều kiện stress như tia UV, nhiệt độ cao, hóa chất gây tổn thương DNA, hoặc thiếu dinh dưỡng, prophage có thể được kích hoạt. Protein repressor bị vô hiệu hóa, dẫn đến việc cắt prophage ra khỏi nhiễm sắc thể vi khuẩn (excision) và chuyển sang chu trình phân giải
\end{enumerate}

Quyết định giữa chu trình phân giải và lysogen phụ thuộc vào nhiều yếu tố, bao gồm trạng thái sinh lý của tế bào chủ, mật độ vi khuẩn, và các tín hiệu môi trường. Cơ chế phân tử điều khiển quyết định này thường liên quan đến sự cạnh tranh giữa các protein điều hòa như CI repressor (ưu tiên lysogen) và Cro protein (ưu tiên phân giải) trong phage lambda.

\subsection{Vai trò sinh thái và ứng dụng}

\subsubsection{Vai trò trong sinh thái vi sinh vật}

Thực khuẩn thể đóng vai trò quan trọng trong sinh thái vi sinh vật:

\begin{itemize}
    \item \textbf{Điều hòa quần thể vi khuẩn:} Phage kiểm soát mật độ và cấu trúc quần thể vi khuẩn thông qua quá trình "kill the winner", trong đó các chủng vi khuẩn chiếm ưu thế bị phage tiêu diệt, tạo cơ hội cho các chủng khác phát triển \cite{koskella2014bacteria}

    \item \textbf{Chuyển gen ngang:} Phage là động lực chính của horizontal gene transfer giữa vi khuẩn thông qua transduction, góp phần vào sự tiến hóa và đa dạng di truyền của vi khuẩn

    \item \textbf{Chu trình dinh dưỡng:} Phage giải phóng các chất dinh dưỡng từ tế bào vi khuẩn bị phân giải, đóng góp vào chu trình carbon và các nguyên tố khác trong hệ sinh thái

    \item \textbf{Lysogenic conversion:} Prophage có thể mang các gen mã hóa yếu tố độc lực hoặc kháng kháng sinh, thay đổi kiểu hình của vi khuẩn chủ
\end{itemize}

\subsubsection{Ứng dụng trong y sinh học}

Việc phân loại chính xác lối sống phage có ý nghĩa quan trọng cho các ứng dụng y sinh học:

\begin{itemize}
    \item \textbf{Liệu pháp phage (Phage therapy):} Sử dụng phage để điều trị nhiễm trùng vi khuẩn, đặc biệt là các vi khuẩn kháng kháng sinh \cite{gorski2016phage,azimi2019phage}. Chỉ có phage độc lực mới phù hợp cho liệu pháp vì chúng tiêu diệt vi khuẩn một cách hiệu quả mà không tích hợp vào bộ gen chủ

    \item \textbf{An toàn sinh học:} Phage ôn hòa có thể truyền các gen độc lực hoặc kháng kháng sinh cho vi khuẩn bệnh nguyên thông qua lysogenic conversion, do đó cần được loại trừ khỏi các ứng dụng trị liệu \cite{cobian2016viruses}

    \item \textbf{Kiểm soát thực phẩm:} Phage được sử dụng để kiểm soát vi khuẩn gây bệnh trong thực phẩm, với yêu cầu nghiêm ngặt về việc sử dụng chỉ phage độc lực

    \item \textbf{Công nghệ sinh học:} Phage được sử dụng trong phage display, vaccine development, và drug delivery, với các yêu cầu khác nhau về lối sống phage tùy thuộc vào ứng dụng cụ thể
\end{itemize}

Do đó, việc phát triển các phương pháp tính toán chính xác và hiệu quả để phân loại lối sống phage từ dữ liệu trình tự gen là cần thiết để hỗ trợ cả nghiên cứu cơ bản về sinh thái vi sinh vật và các ứng dụng thực tiễn trong y học và công nghệ sinh học.
